\documentclass[a4paper, 11pt]{article}
\usepackage[utf8]{inputenc}
\usepackage[T1]{fontenc}
\usepackage{lmodern}
\usepackage[margin=1in]{geometry}
\usepackage{amsmath, amssymb}
\usepackage{booktabs}
\usepackage[protrusion=true,expansion=false]{microtype}
\usepackage{xcolor}
\usepackage{listings}
\usepackage{tcolorbox}
\usepackage{url}
\usepackage{graphicx}

\definecolor{codegreen}{rgb}{0.15, 0.55, 0.15}
\definecolor{codegray}{rgb}{0.45, 0.45, 0.45}
\definecolor{codeblue}{rgb}{0.12, 0.35, 0.75}
\definecolor{backcolour}{rgb}{0.96, 0.97, 0.98}
\definecolor{termbg}{rgb}{0.14, 0.15, 0.17}
\definecolor{termfg}{rgb}{0.92, 0.92, 0.92}

\lstdefinestyle{erlstyle}{
    backgroundcolor=\color{backcolour},   
    commentstyle=\color{codegreen}\itshape,
    keywordstyle=\color{codeblue}\bfseries,
    numberstyle=\tiny\color{codegray},
    stringstyle=\color{codegray},
    basicstyle=\ttfamily\footnotesize,
    breaklines=true,                 
    keepspaces=true,                 
    tabsize=2,
    frame=single,
    rulecolor=\color{black!10}
}
\lstset{style=erlstyle}

\tcbset{
    terminalbox/.style={
        colback=termbg,
        colupper=termfg,
        colframe=black!65,
        arc=3pt,
        boxrule=1pt,
        fontupper=\ttfamily\footnotesize,
        left=6pt, right=6pt, top=5pt, bottom=5pt
    }
}

\title{\textbf{Groupy}}
\author{\textbf{Rodrigo Perestrelo, Lennard Last} \\ ID2201 Distributed Systems, Basic Course}
\date{\today{}}

\begin{document}

\maketitle

\section{Introduction}

The objective of this assignment is to design, implement, and empirically evaluate a group membership service that guarantees atomic multicast within view synchrony. The service coordinates the state of multiple distributed application-layer processes, ensuring they execute an identical sequence of state changes (represented visually through cyclic RGB color rotations).

\section{Iterative Architectural Progression}

We developed the system through three distinct iterations, progressively addressing failure detection and atomic delivery across node crashes.

\subsection{Iteration 1: Centralized Multicast with Dynamic Joins (GMS1)}
The initial architecture, \texttt{gms1}, provides basic group membership and multicast serialization under the assumption of zero crashes. A single node is designated as the group leader, while all other participants act as slaves:
\begin{itemize}
    \item \textbf{Sequenced Multicast:} All application state changes are submitted via \texttt{\{mcast, Msg\}}. Slaves forward requests to the leader, which serializes them through \texttt{bcast/3} across the group, enforcing total ordering.
    \item \textbf{Deterministic Admission:} New nodes contact an existing peer with \texttt{\{join, Wrk, Peer\}}. The request reaches the leader, which appends the new peer to the tail of \texttt{Slaves} and \texttt{Group}, multicasting an updated view.
    \item \textbf{State Catch-Up on Admission:} When admitted, a joining node multicasts a \texttt{\{state\_request, Ref\}}. Because all multicasts are totally ordered, existing nodes observe this request at the same logical index and reply with their current state (\texttt{\{state, Ref, Color\}}), allowing the new node to synchronize before applying deferred pipeline changes.
\end{itemize}

\noindent \textbf{The Fatal Flaw:} The leader represents an unmitigated single point of failure. If the leader halts, slaves block permanently in their forwarding loop.

\subsection{Iteration 2: Failure Detection and Leader Election (GMS2)}
To eliminate deadlocks on leader crash, \texttt{gms2} introduces asynchronous failure detection and autonomous succession:
\begin{itemize}
    \item \textbf{Targeted Process Monitors:} Slaves monitor exclusively the current leader via \texttt{erlang:monitor(process, Leader)}.
    \item \textbf{Deterministic Election:} Upon receiving a \texttt{'DOWN'} signal, slaves enter \texttt{election/4}. Because all nodes preserve identical list ordering, the first element (\texttt{[Self | Rest]}) autonomously assumes the leadership role and broadcasts the updated view. Remaining peers redirect their monitor to the new leader.
    \item \textbf{Join Timeouts:} A \texttt{?timeout} clause prevents joining nodes from hanging indefinitely if a leader crash coincides with a join request.
\end{itemize}

\begin{lstlisting}[language=erlang]
election(Id, Master, Slaves, Group) ->
    Self = self(),
    case Slaves of
        [Self | Rest] ->
            bcast(Id, {view, Slaves, Group}, Rest),
            Master ! {view, Group},
            leader(Id, Master, Rest, Group);
        [Leader | Rest] ->
            erlang:monitor(process, Leader),
            slave(Id, Master, Leader, Rest, Group)
    end.
\end{lstlisting}

\subsection{Iteration 3: Atomic Multicast with View Synchrony (GMS3)}
While \texttt{gms2} survives clean exits, it violates atomic delivery when a leader crashes mid-broadcast. \texttt{gms3} enforces view synchrony by introducing message sequence numbers $N$, caching the \texttt{Last} delivered broadcast, and leveraging two structural invariants:
\begin{enumerate}
    \item \textbf{Sequential Dispatch Invariant:} The leader transmits messages strictly in the order of the \texttt{Slaves} list.
    \item \textbf{Succession Invariant:} The first slave in \texttt{Slaves} is always elected as the next leader.
\end{enumerate}

\noindent \textbf{Deduction:} If the dead leader delivered a message to \emph{any} node in the group before crashing, the designated successor (\texttt{Self}) is guaranteed to have received it. Upon election, the new leader resends \texttt{Last} to all surviving peers before broadcasting the newly installed view:

\begin{lstlisting}[language=erlang]
election(Id, Master, N, Last, Slaves, Group) ->
    Self = self(),
    case Slaves of
        [Self | Rest] ->
            bcast(Id, Last, Rest),
            bcast(Id, {view, N, Slaves, Group}, Rest),
            Master ! {view, Group},
            leader(Id, Master, N + 1, Rest, Group);
        [NewLeader | Rest] ->
            erlang:monitor(process, NewLeader),
            slave(Id, Master, NewLeader, N, Last, Rest, Group)
    end.
\end{lstlisting}

\noindent \textbf{Duplicate Suppression:} To prevent peers that already received \texttt{Last} prior to the crash from processing it twice (which would rotate state out of phase), slaves filter out duplicate sequence numbers ($I < N$):

\begin{lstlisting}[language=erlang]
{msg, I, _} when I < N ->
    slave(Id, Master, Leader, N, Last, Slaves, Group);
{msg, N, Msg} ->
    Master ! Msg,
    Last2 = {msg, N, Msg},
    slave(Id, Master, Leader, N + 1, Last2, Slaves, Group);
{view, I, _, _} when I < N ->
    slave(Id, Master, Leader, N, Last, Slaves, Group);
{view, N, [Leader | Slaves2], Group2} ->
    Master ! {view, Group2},
    Last2 = {view, N, [Leader | Slaves2], Group2},
    slave(Id, Master, Leader, N + 1, Last2, Slaves2, Group2);
\end{lstlisting}

\section{Empirical Evaluation}

\subsection{Evaluation 1: GMS1 Failure-Free Coordination and Leader Deadlock}
We deployed an initial 3-node cluster under \texttt{gms1} with a cycle interval of 1000 ms to evaluate behavior under both static coordination and leader termination:
\begin{tcolorbox}[terminalbox]
1> c(gui), c(worker), c(test), c(gms1). \\
\{ok,gms1\} \\
2> W1 = test:first(1, gms1, 1000). \\
<0.92.0> \\
3> W2 = test:add(2, gms1, W1, 1000). \\
<0.95.0> \\
4> W3 = test:add(3, gms1, W1, 1000). \\
<0.98.0> \\
5> exit(W1, kill). \\
true
\end{tcolorbox}
\textbf{Observed Result:} 
\begin{itemize}
    \item \textbf{Failure-Free Operation:} All three nodes updated their states simultaneously, confirming that centralized sequencer routing over FIFO links satisfies total ordering when nodes do not crash.
    \item \textbf{Leader Crash Deadlock:} Upon executing \texttt{exit(W1, kill)}, surviving workers \texttt{W2} and \texttt{W3} completely stopped updating and froze permanently. Because \texttt{gms1} lacks process monitors and election mechanisms, surviving slaves continued forwarding state requests to the deceased leader PID without error recovery, empirically demonstrating the single-point-of-failure limitation.
\end{itemize}

\subsection{Evaluation 2: GMS2 State Divergence Under Crash Injection}
We initialized a 4-node cluster under \texttt{gms2} with crash injection probability $\text{arghh} = 150$:
\begin{tcolorbox}[terminalbox]
1> c(gms2). \\
\{ok,gms2\} \\
2> W1 = test:first(1, gms2, 500). \\
3> W2 = test:add(2, gms2, W1, 500). \\
4> W3 = test:add(3, gms2, W1, 500). \\
5> W4 = test:add(4, gms2, W1, 500). \\
leader 1: crash
\end{tcolorbox}
\textbf{Observed Result:} Worker 1 crashed while multicasting a color update. Worker 2 received the message, updated its color, and promoted itself to leader upon receiving the monitor signal. However, Workers 3 and 4 missed the message prior to the crash. Consequently, Worker 2 displayed a color state completely out of phase with Workers 3 and 4. As new updates were broadcast, this offset persisted permanently, confirming that \texttt{gms2} fails to preserve consistency under mid-broadcast failure.

\subsection{Evaluation 3: GMS3 View Synchrony and Self-Healing}
We repeated the crash experiment using \texttt{gms3}:
\begin{tcolorbox}[terminalbox]
1> c(gms3). \\
\{ok,gms3\} \\
2> W1 = test:first(1, gms3, 500). \\
3> W2 = test:add(2, gms3, W1, 500). \\
4> W3 = test:add(3, gms3, W1, 500). \\
5> W4 = test:add(4, gms3, W1, 500). \\
leader 1: crash
\end{tcolorbox}
\textbf{Observed Result:} Worker 1 crashed mid-broadcast. Worker 2 detected the crash, retransmitted \texttt{Last} to Workers 3 and 4, and broadcast the new view. Workers 3 and 4 caught up on the missing transition, while any duplicate on Worker 2 was dropped by the $I < N$ guard. All surviving windows continued rotating colors in lockstep.

\subsection{Evaluation 4: The Rolling Group Experiment}
To answer the core question of Section 3.4 (``\emph{Can we keep a group rolling by adding more nodes as existing nodes die?}''), we subjected \texttt{gms3} to continuous leader crashes while dynamically adding new nodes:
\begin{tcolorbox}[terminalbox]
leader 1: crash \\
6> W5 = test:add(5, gms3, W2, 500). \\
leader 2: crash \\
7> W6 = test:add(6, gms3, W3, 500).
\end{tcolorbox}
\textbf{Observed Result:} The group remained operational indefinitely. Whenever a leader crashed, the next sequential slave assumed leadership without state divergence. Newly admitted workers requested the current color, caught up with in-flight messages, and seamlessly merged into the rolling view.

\section{Optional Task for Extra Bonus: What Could Possibly Go Wrong}

\noindent \textit{[Placeholder.]}

\section{Conclusions}

This laboratory demonstrated the design of an atomic multicast protocol operating under view synchrony. Centralized FIFO multicasting (\texttt{gms1}) provides total ordering in static topologies but deadlocks upon leader crash. Adding naive failure detection (\texttt{gms2}) enables failover but violates atomic delivery during mid-transmission failures. By combining deterministic list ordering, sequential succession, retransmission of the last delivered message, and monotonic sequence filtering, \texttt{gms3} guarantees that all correct processes agree on an identical state history across dynamic views.

\end{document}